\subsection{Comparison to Related Work}
\label{performance-estimates}


For comparison,
we next expand on our performance estimates for two related schemes given in Table~\ref{table:efficiency}.
We do not give exact computation estimates,
but instead give estimates for computation that dominate.
For bandwidth,
we estimate for a $256$-bit elliptic curve,
and so estimate scalers as $32$ bytes and group elements as $33$ bytes.

\paragraph{MuSig-DN~\cite{NickRSW20}.}
Nick et al.~\cite{NickRSW20} present the proof size
as $1124$~bytes, which agrees with our measurements.
Further, each participant sends one group element in round one,
and outputs one field element in round three, yielding
the total of $1189$ bytes we report in Table~\ref{table:efficiency}.

While Nick et al.\ give total computational overhead for MuSig-DN for a $256$-bit curve,
we give more detail as to computation that dominates when performing signing operations.

The computational overhead of MuSig-DN is dominated by the cost of generating and verifying Bulletproof~\cite{BunzBBPWM18} proofs.
For each signing operation,
each participant must generate one proof and verify $(n-1)$ other proofs.
While the computational overhead of Bulletproofs scales linearly relative to the size of the circuit,
batch verification and multi-scalar multiplication allow for improved performance.

Proofs in MuSig-DN require a circuit of 2030 multiplication gates~\cite{NickRSW20},
and estimates are given with respect to a $256$-bit elliptic curve group.
Naively, this overhead would translate into roughly $2030n$ operations for each signer,
due to the requirement that each signer generates a proof and verifies $n-1$ other proofs.

However, when considering speedups offered by multi-scalar multiplications and batch verification,
the overhead for a signer becomes roughly $2030$ operations,
plus additional (small) terms linear in $n$.
When generating a MuSig-DN proof using Bulletproofs,
we estimate that the dominant computational overhead is roughly $6 \cdot 2030$ multi-scalar multiplications,
to commit to the circuit constraints and additional factors.
For batch verifying $n-1$ MuSig-DN proofs,
we estimate the dominant computational overhead to the verifier is $2030 + 24(n-1)$ scalar multiplications.


\paragraph{GKMN21~\cite{GarillotKMN21}.}
Garillot et al.~\cite{GarillotKMN21} give computational and bandwidth estimates for a single sender and receiver,
but each player will be involved in $t-1$ concurrent sending and receiving sessions with all other players.
As such, the estimate we give in Table~\ref{table:efficiency} is for the total sending and receiving overhead for each signer in a signing session with $t$ participants.

\paragraph{Arctic.}
For bandwidth,
each party sends one group element and one field element as the output to the first signing query,
and then outputs a single field element in the second signing round.

For computation,
each party performs $\DVRFOnem.\generate$ once for each signing round,
resulting in ${n-1 \choose t-1}$ field and hash operations, and one
group multiplication, for each execution.
In the second round of signing,
each party performs
$\DVRFOnem.\generate$ again, for another
${n-1 \choose t-1}$ field and hash operations, and one
group multiplication, and also
$\DVRFOnem.\verify$ and $\DVRFOnem.\combinepub$,
which require $2t^2 -t$ group multiplications.
